is also representable by a closed immersion.}
\item \emph{@@M0@@ is of finite type over @@M1@@, that is, the family of classes of invertible sheaves @@M2@@-equivalent to @@M3@@ on the fibers of @@M4@@ is bounded.}
\end{enumerate}

Indeed, in (i) (resp. (ii)), we must treat the case of a morphism @@M5@@, where @@M6@@ is an arbitrary @@M7@@-scheme, and show that the inverse image of @@M8@@ is representable by an open (resp. and closed) subscheme of @@M9@@. It suffices to take @@M10@@ over an affine open @@M11@@ of @@M12@@. Then @@M13@@ is the limit of the filtering inverse system of affine @@M14@@-schemes @@M15@@, where @@M16@@ runs through the finite-type @@M17@@-subalgebras of @@M18@@. By 3.1, the morphism in question comes from a morphism @@M19@@. We may therefore suppose that @@M20@@ is of finite type over @@M21@@.

The morphism @@M22@@ corresponds to an f.p.p.f. morphism @@M23@@ and an invertible sheaf @@M24@@ on @@M25@@. By f.p.p.f. descent, we may replace @@M26@@ by @@M27@@. After the base change @@M28@@, taking into account the compatibility of the formation of @@M29@@ with every base change (4.1), we are thus reduced in (i) and (ii) to studying the case of a section @@M30@@, defined by an invertible sheaf @@M31@@ on @@M32@@.

We must prove that the set @@M33@@ of @@M34@@ such that @@M35@@ is @@M36@@-equivalent to @@M37@@ is open, and also closed under the conditions of (ii); since @@M38@@ is noetherian, this means three things: for an @@M39@@,
\begin{enumerate}[label=(\alph*)]
\item every generalization of @@M40@@ remains in @@M41@@.
\item every specialization of @@M42@@, near @@M43@@, remains in @@M44@@.
\item under the conditions of (ii), every specialization of @@M45@@ remains in @@M46@@.
\end{enumerate}

